\section{Results}\label{sec:results}
The main result is \S\ref{sec:main}: DA-v2 Metric-Indoor fine-tuned with the iPad's own LiDAR as the teacher (R3)
against the same checkpoint pretrained, both measured against the Faro laser scanner---the raw ground truth that
ARKitScenes ships. \S\ref{sec:finetune} ablates the teacher; \S\ref{sec:exp1}--\ref{sec:exp5} (Exp.~1--5) price
every other depth source through the same pipeline and serve as context for where the remaining error comes from.
3D numbers are copied from \path{experiments/results/<exp>/table.md} (mean $\pm$ std over scenes) and
\code{per\_scene.md}, generated by \code{roomscan report}; macros are generated by \path{scripts/paper_numbers.py}.
Every 2D number uses protocol~2 (pixels selected by the ground truth, prediction clipped; \S\ref{sec:metrics}).
Units are cm; $F$ is the F-score at $\tau=$\SI{5}{cm}.

\subsection{Main result: LiDAR-teacher fine-tune (R3) vs.\ the pretrained model}\label{sec:main}
\textbf{Question.} If a metric monocular model is taught with depth from the iPad's own LiDAR---no laser scanner in
the training loop---is its depth really better than the same model pretrained, \emph{when measured against the laser
scanner}? The two rows differ only in the weights of the DPT neck and metric head: \emph{pretrained} is
\code{Depth-Anything-V2-Metric-Indoor-Large} (trained on synthetic Hypersim) and \emph{R3} is that checkpoint
fine-tuned on 24 ARKitScenes rooms with \code{lowres\_depth} labels (\S\ref{sec:ftsetup}). At test time both see the
same \code{vga\_wide} frames, the same upright rotation and the same $518\times686$ input, with no alignment of any kind.
No room in this subsection was used for training, checkpoint selection or any hyper-parameter (split by
\code{visit\_id}). We measure at two levels: \emph{3D}, the mesh against the reference fused from Faro depth, on the six
test rooms; and \emph{2D}, per-frame depth against the dataset's Faro \code{highres\_depth} directly, on those six rooms
\emph{and on \rsVN{} further rooms from the ARKitScenes Validation fold} (Exp.~7, \rsVFRAMES{} Faro frames).

\begin{table*}[t]
\centering\small
\caption{Main result on the six test rooms: RGB only, no alignment, measured against Faro. The two reference rows come
from Exp.~1 and need a sensor at run time.}
\label{tab:main}
\begin{adjustbox}{max width=\textwidth}
\begin{tabular}{@{}llrrrrrrrr@{}}
\toprule
Model & Teacher & Chamfer & Acc. & Compl. & \textbf{F@5} & R@10 & \textbf{AbsRel} & $\delta_1$ & RMSE \\
\midrule
DA-v2 Metric-Indoor, pretrained & --- (Hypersim) & $\rsMETRICCM \pm \rsMETRICSD$ & \rsMETRICACC & \rsMETRICCOMP & \rsMETRICF & \rsMETRICRten & \rsMETRICSIXABSREL & \rsMETRICSIXDELTAONE & \rsMETRICSIXRMSE \\
$+$ fine-tune \textbf{R3} & iPad LiDAR, 24 rooms & $\mathbf{\rsFTLIDARALLCM \pm \rsFTLIDARALLSD}$ & \textbf{\rsFTLIDARALLACC} & \textbf{\rsFTLIDARALLCOMP} & \textbf{\rsFTLIDARALLF} & \textbf{\rsFTLIDARALLRten} & \textbf{\rsFTLIDARALLABSREL} & \textbf{\rsFTLIDARALLDELTAONE} & \textbf{\rsFTLIDARALLRMSE} \\
\midrule
\emph{ref.:} iPad LiDAR (sensor) & --- & $\rsLIDARCM \pm \rsLIDARSD$ & \rsLIDARACC & \rsLIDARCOMP & \rsLIDARF & \rsLIDARRten & \rsLIDARABSREL & \rsLIDARDELTAONE & \rsLIDARRMSE \\
\emph{ref.:} DA-v2 L $+$ 200 sparse pts/frame & --- & $\rsSPARSECM \pm \rsSPARSESD$ & \rsSPARSEACC & \rsSPARSECOMP & \rsSPARSEF & \rsSPARSERten & \rsSPARSEABSREL & \rsSPARSEDELTAONE & \rsSPARSERMSE \\
\bottomrule
\end{tabular}
\end{adjustbox}
\end{table*}

\begin{table}[t]
\centering\small
\caption{Main result per test room (order of Table~\ref{tab:scenes}): Chamfer (cm) and AbsRel.}
\label{tab:mainscene}
\begin{adjustbox}{max width=\columnwidth}
\begin{tabular}{@{}lrrrrrrr@{}}
\toprule
 & 4244 & 4733 & 4711 & 4289 & 4742 & 4526 & mean \\
\midrule
Chamfer, pretrained & \rsMETRICa & \rsMETRICb & \rsMETRICc & \rsMETRICd & \rsMETRICe & \rsMETRICf & \rsMETRICCM \\
Chamfer, R3 & \rsFTLIDARALLa & \rsFTLIDARALLb & \rsFTLIDARALLc & \rsFTLIDARALLd & \rsFTLIDARALLe & \rsFTLIDARALLf & \rsFTLIDARALLCM \\
AbsRel, pretrained & \rsMETRICARa & \rsMETRICARb & \rsMETRICARc & \rsMETRICARd & \rsMETRICARe & \rsMETRICARf & \rsMETRICSIXABSREL \\
AbsRel, R3 & \rsFTALLARa & \rsFTALLARb & \rsFTALLARc & \rsFTALLARd & \rsFTALLARe & \rsFTALLARf & \rsFTLIDARALLABSREL \\
\bottomrule
\end{tabular}
\end{adjustbox}
\end{table}

\begin{table}[t]
\centering\small
\caption{Exp.~7: 2D only, \rsVN{} Validation-fold rooms never used for anything else (every second Faro frame,
\rsVFRAMES{} frames); mean $\pm$ std over rooms. Per-room values: \code{experiments/results/exp7\_valfold\_2d/}.}
\label{tab:exp7}
\begin{adjustbox}{max width=\columnwidth}
\begin{tabular}{@{}lrrr@{}}
\toprule
Row & AbsRel & $\delta_1$ & RMSE (cm) \\
\midrule
DA-v2 Metric-Indoor, pretrained & $\rsVMETRICABSREL \pm \rsVMETRICABSRELSD$ & $\rsVMETRICDELTAONE \pm \rsVMETRICDELTAONESD$ & \rsVMETRICRMSE \\
$+$ fine-tune \textbf{R3} & $\mathbf{\rsVFTALLABSREL \pm \rsVFTALLABSRELSD}$ & $\mathbf{\rsVFTALLDELTAONE \pm \rsVFTALLDELTAONESD}$ & \textbf{\rsVFTALLRMSE} \\
\emph{ref.:} iPad LiDAR (sensor) & $\rsVLIDARABSREL \pm \rsVLIDARABSRELSD$ & $\rsVLIDARDELTAONE \pm \rsVLIDARDELTAONESD$ & \rsVLIDARRMSE \\
\bottomrule
\end{tabular}
\end{adjustbox}
\end{table}

\begin{figure}[t]
\centering
\includegraphics[width=\columnwidth]{finetune_absrel}
\caption{Per-room AbsRel against Faro depth: pretrained (x) vs.\ the LiDAR-teacher fine-tune R3 (y), for the six test
rooms (Exp.~6) and the \rsVN{} Validation-fold rooms (Exp.~7). Below the diagonal $=$ the fine-tune is better.}
\label{fig:finetuneabsrel}
\end{figure}

\begin{enumerate}
\item \textbf{The LiDAR-teacher fine-tune beats the pretrained model in every test room.} Chamfer
\SI{\rsMETRICCM}{cm}$\rightarrow$\SI{\rsFTLIDARALLCM}{cm} ($\rsMAINGAIN\times$; \SI{\rsMAINDIFFCM \pm \rsMAINDIFFSD}{cm}
better per room, \rsMAINWINS/\rsMAINN{} rooms, exact Wilcoxon $p=\rsMAINP$, the smallest attainable at $n=6$) and
\fscore{} \rsMETRICF$\rightarrow$\rsFTLIDARALLF{} (Tables~\ref{tab:main}, \ref{tab:mainscene}). The room where the
pretrained model fails worst (kitchen 45261556) is still R3's worst, but three times better.
\item \textbf{The 3D error falls because misplaced surfaces fall.} Accuracy (pred$\rightarrow$ref) drops
\SI{\rsMETRICACC}{cm}$\rightarrow$\SI{\rsFTLIDARALLACC}{cm}, more than completeness
(\SI{\rsMETRICCOMP}{cm}$\rightarrow$\SI{\rsFTLIDARALLCOMP}{cm}): the pretrained model builds the room but too far away
(it over-estimates distance by 15--31\,\%, \S\ref{sec:metric}); R3 brings the surfaces back near where they belong.
\item \textbf{The same holds in 2D on rooms that took part in nothing.} On the \rsVN{} Validation-fold rooms AbsRel
drops \rsVMETRICABSREL$\rightarrow$\rsVFTALLABSREL{} ($\rsVGAIN\times$; R3 better in \rsVWINS/\rsVN{} rooms, exact
Wilcoxon $p=\rsVP$) and $\delta_1$ rises \rsVMETRICDELTAONE$\rightarrow$\rsVFTALLDELTAONE{} (better in
\rsVWINSDELTA/\rsVN) (Table~\ref{tab:exp7}, Fig.~\ref{fig:finetuneabsrel}). The one room that does not improve is \rsVEXCEPT{}
(AbsRel \rsVEXCEPTPRE$\rightarrow$\rsVEXCEPTFT, $\delta_1$ \rsVEXCEPTDPRE$\rightarrow$\rsVEXCEPTDFT)---one of the
pretrained model's best rooms and R3's worst; in the other 19, R3's AbsRel is 21--64\,\% of the pretrained one
(median ratio over all \rsVN{} rooms \rsVMEDRATIO). The improvement is not a property of the six test rooms chosen at the start of the
project.
\item \textbf{It does not reach a sensor.} R3 is still $\sim$\SI{13}{cm} behind the iPad's LiDAR in Chamfer and
\SI{\rsFTALLGAPSPARSE}{cm} behind a per-frame fit on 200 sparse points; in 2D on the Validation fold LiDAR has AbsRel
\rsVLIDARABSREL{} vs.\ R3's \rsVFTALLABSREL. Fine-tuning fixes the scale of the \emph{domain}; it does not fix the
scale of \emph{the frame being looked at} (\S\ref{sec:whyft}). For a device with neither LiDAR nor trustworthy VIO
points, R3 is the only row in this paper that works without calibration at run time.
\end{enumerate}

\begin{figure*}[t]
\centering
\includegraphics[width=\textwidth]{exp6_finetune_topdown_grid}
\caption{Exp.~6 top-down error maps (rows $=$ scenes, columns $=$ the three fine-tunes and the pretrained metric
model, colour $=$ distance to the reference, 0--\SI{10}{cm}). The pretrained column is red almost everywhere---the
room is built, at the wrong distance. After fine-tuning the walls and floor come back (blue/cyan) and the error
concentrates at object edges and far walls; where a fine-tune is weaker (45261556 and 47333774 with the LiDAR
teacher) it shows as a thick red rim around the room rather than a displaced room.}
\label{fig:exp6}
\end{figure*}

\subsection{Teacher ablation and a Level-0 baseline}\label{sec:finetune}
R3 is one of three fine-tunes with the same recipe (frozen DINOv2 backbone, DPT neck $+$ head, SiLog loss) that differ
in label and number of rooms: Faro \code{highres\_depth} on 10 rooms (R1, 3{,}594 frames), the iPad's LiDAR on the
same 10 rooms (R2, 8{,}656 frames), and LiDAR on 24 rooms (R3, the 10 $+$ 14 rooms with no laser scan at all,
19{,}896 frames). R1 and R2 share rooms but \emph{not frames}: Faro exists only on the frames Apple kept
($\sim$4~fps) whereas LiDAR accompanies every colour frame (every third taken, $\sim$10~fps), so at the same 6{,}000 steps
R2 sees $\sim$2.4$\times$ more distinct views than R1. The R1--R2 comparison therefore changes both the label and view
diversity. As a Level-0 reference the table also carries Depth Pro~\cite{depthpro}, a metric model of the other family
(\S\ref{sec:related}): no fine-tuning, and \emph{handed the camera's true focal length} (\SI{533}{px}) rather than the
one it predicts, which is the strongest form of that baseline.

\begin{table}[h]
\centering\small
\caption{Exp.~6, all rows: same pipeline, same six scenes, RGB-only, no alignment. Teacher $=$ the depth used as the training
label; none of these rows uses a depth sensor at test time.}
\label{tab:exp6}
\begin{adjustbox}{max width=\columnwidth}
\begin{tabular}{@{}llrrrr@{}}
\toprule
Model & Teacher & Chamfer & \textbf{F@5} & AbsRel & $\delta_1$ \\
\midrule
Depth Pro (intrinsics given) & --- & $\rsDEPTHPROCM \pm \rsDEPTHPROSD$ & \rsDEPTHPROF & \rsDEPTHPROABSREL & \rsDEPTHPRODELTAONE \\
Metric-Indoor (pretrained) & --- & $\rsMETRICCM \pm \rsMETRICSD$ & \rsMETRICF & \rsMETRICSIXABSREL & \rsMETRICSIXDELTAONE \\
$+$ fine-tune (R1) & Faro & $\mathbf{\rsFTFAROCM \pm \rsFTFAROSD}$ & \textbf{\rsFTFAROF} & \rsFTFAROABSREL & \rsFTFARODELTAONE \\
$+$ fine-tune (R2) & ARKit LiDAR & $\rsFTLIDARCM \pm \rsFTLIDARSD$ & \rsFTLIDARF & \rsFTLIDARABSREL & \textbf{\rsFTLIDARDELTAONE} \\
$+$ fine-tune (R3) & LiDAR, 24 scenes & $\rsFTLIDARALLCM \pm \rsFTLIDARALLSD$ & \rsFTLIDARALLF & \textbf{\rsFTLIDARALLABSREL} & \rsFTLIDARALLDELTAONE \\
\bottomrule
\end{tabular}
\end{adjustbox}
\end{table}

\begin{table}[h]
\centering\small
\caption{Exp.~6 Chamfer (cm) per scene; column order as in Table~\ref{tab:exp1scene}.}
\label{tab:exp6scene}
\begin{adjustbox}{max width=\columnwidth}
\begin{tabular}{@{}lrrrrrrr@{}}
\toprule
Teacher & 4244 & 4733 & 4711 & 4289 & 4742 & 4526 & mean \\
\midrule
Depth Pro (no fine-tune) & \rsDEPTHPROa & \rsDEPTHPROb & \rsDEPTHPROc & \rsDEPTHPROd & \rsDEPTHPROe & \rsDEPTHPROf & \rsDEPTHPROCM \\
none (pretrained) & \rsMETRICa & \rsMETRICb & \rsMETRICc & \rsMETRICd & \rsMETRICe & \rsMETRICf & \rsMETRICCM \\
Faro (R1) & \rsFTFAROa & \rsFTFAROb & \rsFTFAROc & \rsFTFAROd & \rsFTFAROe & \rsFTFAROf & \rsFTFAROCM \\
LiDAR (R2) & \rsFTLIDARa & \rsFTLIDARb & \rsFTLIDARc & \rsFTLIDARd & \rsFTLIDARe & \rsFTLIDARf & \rsFTLIDARCM \\
LiDAR $\times$24 (R3) & \rsFTLIDARALLa & \rsFTLIDARALLb & \rsFTLIDARALLc & \rsFTLIDARALLd & \rsFTLIDARALLe & \rsFTLIDARALLf & \rsFTLIDARALLCM \\
\bottomrule
\end{tabular}
\end{adjustbox}
\end{table}

\begin{enumerate}
\item \textbf{The laser teacher is not measurably better than the LiDAR teacher.} R1 \SI{\rsFTFAROCM}{cm}, R3
\SI{\rsFTLIDARALLCM}{cm} and R2 \SI{\rsFTLIDARCM}{cm} all lie within the across-room std (4.0--\SI{4.7}{cm}), and no
row wins every room (R1 beats R3 in three rooms, R3 wins two, one tie; Table~\ref{tab:exp6scene}). Given the frame
difference above, the honest claim is that a phone's depth sensor can replace the laser scanner as the teacher at no
measurable cost on these six rooms---not that the two are equal, or that either is better.
\item \textbf{Fourteen extra laser-free rooms do not help measurably} (R3 \SI{\rsFTLIDARALLCM}{cm} vs.\ R2
\SI{\rsFTLIDARCM}{cm}). R3 was the best run on the validation fold (AbsRel 0.133 vs.\ 0.142, measured during training under protocol~1) and is not the best here:
three validation rooms are too few to rank checkpoints this close (\S\ref{sec:limitations}). We take R3 as the main
model for reasons that do not depend on the test numbers: it is the only recipe that scales without a laser scanner,
and it was the best run on validation.
\item \textbf{Knowing the camera does not rescue metric depth.} Depth Pro, given the true intrinsics, is
\SI{\rsDEPTHPROCM}{cm} with \fscore{} \rsDEPTHPROF---worse than the pretrained DA-v2 Metric row it was meant to
improve on, and $\sim$10$\times$ the fine-tuned rows. Its structure is right (the surfaces of a frame land in the
right places relative to each other) and its distances are not: it predicts a \SI{48}{\degree} field of view where
the camera has \SI{62}{\degree}, and even after removing all scale per frame its AbsRel is 0.24--0.54 against
0.04--0.13 for DA-v2 Metric on the same frames. The failure is specific to this input---VGA
($640\times480$), close indoor range, motion blur---and we verified the wrapper against a synthetic room with exact
ground truth (\fscore{}-level agreement, AbsRel 0.008 after one scale), so it is the model on this data, not the
plumbing. The practical reading for a product team: a model that takes intrinsics is not a substitute for depth
labels from the target device.
\end{enumerate}

\subsection{Context, Exp.~1: the price of each depth source}\label{sec:exp1}
\begin{table*}[t]
\centering\small
\caption{Same pipeline, \SI{4}{cm} voxels, every Faro frame; only \code{depth.source} / \code{depth.aligner} change.
\code{sparse} $=$ per-frame fit on 200 LiDAR-sampled points (VIO proxy; upper bound of real VIO).}
\label{tab:exp1}
\begin{adjustbox}{max width=\textwidth}
\begin{tabular}{@{}llrrrrrrrr@{}}
\toprule
Depth source & Aligner & Chamfer & Acc. & Compl. & F@2 & \textbf{F@5} & F@10 & AbsRel (2D) & Time/scan (s) \\
\midrule
Faro laser (\code{gt}) & identity & $\mathbf{\rsGTCM \pm \rsGTSD}$ & \rsGTACC & \rsGTCOMP & \rsGTFtwo & \textbf{\rsGTF} & \rsGTFten & 0 & 1.6 \\
ARKit LiDAR (\code{lidar}) & identity & $\mathbf{\rsLIDARCM \pm \rsLIDARSD}$ & \rsLIDARACC & \rsLIDARCOMP & \rsLIDARFtwo & \textbf{\rsLIDARF} & \rsLIDARFten & \rsLIDARABSREL & 1.7 \\
DA-v2 L (\code{mono}) & \code{oracle\_frame} & $\mathbf{\rsORACLECM \pm \rsORACLESD}$ & \rsORACLEACC & \rsORACLECOMP & \rsORACLEFtwo & \textbf{\rsORACLEF} & \rsORACLEFten & \rsORACLEABSREL & 366 \\
DA-v2 L (\code{mono}) & \code{sparse\_points} & $\mathbf{\rsSPARSECM \pm \rsSPARSESD}$ & \rsSPARSEACC & \rsSPARSECOMP & \rsSPARSEFtwo & \textbf{\rsSPARSEF} & \rsSPARSEFten & \rsSPARSEABSREL & 360 \\
DA-v2 L (\code{mono}) & \code{per\_scene} & $\mathbf{\rsSCENECM \pm \rsSCENESD}$ & \rsSCENEACC & \rsSCENECOMP & \rsSCENEFtwo & \textbf{\rsSCENEF} & \rsSCENEFten & \rsSCENEABSREL & 351 \\
DA-v2 Metric-Indoor & identity & $\mathbf{\rsMETRICCM \pm \rsMETRICSD}$ & \rsMETRICACC & \rsMETRICCOMP & \rsMETRICFtwo & \textbf{\rsMETRICF} & \rsMETRICFten & \rsMETRICABSREL & 352 \\
\bottomrule
\end{tabular}
\end{adjustbox}
\end{table*}

\begin{table*}[t]
\centering\small
\caption{Chamfer (cm) per scene; column order as in Table~\ref{tab:scenes}. 42444474 is the development scene.}
\label{tab:exp1scene}
\begin{adjustbox}{max width=\textwidth}
\begin{tabular}{@{}lrrrrrrr@{}}
\toprule
Run & 42444474 kitchen/glass & 47333774 bedroom & 47115299 living & 42897743 bathroom & 47429736 bedroom$+$closets & 45261556 kitchen & mean \\
\midrule
gt & \rsGTa & \rsGTb & \rsGTc & \rsGTd & \rsGTe & \rsGTf & \rsGTCM \\
lidar & \rsLIDARa & \rsLIDARb & \rsLIDARc & \rsLIDARd & \rsLIDARe & \rsLIDARf & \rsLIDARCM \\
mono $+$ oracle & \rsORACLEa & \rsORACLEb & \rsORACLEc & \rsORACLEd & \rsORACLEe & \rsORACLEf & \rsORACLECM \\
mono $+$ sparse & \rsSPARSEa & \rsSPARSEb & \rsSPARSEc & \rsSPARSEd & \rsSPARSEe & \rsSPARSEf & \rsSPARSECM \\
mono $+$ per\_scene & \rsSCENEa & \rsSCENEb & \rsSCENEc & \rsSCENEd & \rsSCENEe & \rsSCENEf & \rsSCENECM \\
mono metric & \rsMETRICa & \rsMETRICb & \rsMETRICc & \rsMETRICd & \rsMETRICe & \rsMETRICf & \rsMETRICCM \\
\bottomrule
\end{tabular}
\end{adjustbox}
\end{table*}

What the table says directly:
\begin{enumerate}
\item \textbf{The pipeline loses $\sim$\SI{2}{cm}; LiDAR adds $\sim$\SI{1}{cm}.} The gap \code{lidar}$-$\code{gt} is
\SI{0.98 \pm 0.25}{cm} (0.66--1.33), the same in every room including the bathroom; \fscore{} 0.93 vs 0.97. The rows
differ clearly only at $\tau=$\SI{2}{cm} (0.54 vs 0.86): iPad LiDAR is wrong at the 2--\SI{4}{cm} level, not \SI{5}{cm}.
\item \textbf{Monocular depth with the right scale on every frame (oracle) is
\SI{\rsORACLEGAPLIDARCM \pm \rsORACLEGAPLIDARSD}{cm} behind LiDAR} (0.9--3.8). Bedroom 47333774 (3.8) and the bathroom
(3.4) are farthest; the development kitchen 42444474 (0.9) and kitchen 45261556 (1.6) closest---more room-dependent
than LiDAR's std suggests (\S\ref{sec:rooms}).
\item \textbf{Two hundred sparse metric points per frame recover the oracle.} \code{sparse\_points} lands at
\SI{\rsSPARSECM}{cm}, \SI{\rsSPARSEGAP}{cm} from the oracle and $\rsSPARSEVSSCENE\times$ better than one-off calibration,
with the same per-scene ordering as the oracle (Table~\ref{tab:exp1scene}). The remaining gap to LiDAR
(\SI{\rsSPARSEGAPLIDAR}{cm}) is shape, not scale.
\item \textbf{One calibration per scene (\code{per\_scene}) is $4.5 \pm 1.1\times$ worse than the oracle} (3.1--6.4$\times$),
and---importantly---\textbf{the absolute values cluster at 20--\SI{26}{cm} in every room} (std \SI{\rsSCENESD}{cm}, lower than
the oracle's): not one room failing, but a property of affine-invariant models with one scale per scene (\S\ref{sec:drift}).
\item \textbf{The metric model is worst (\SI{\rsMETRICCM}{cm}, $F$ \rsMETRICF) and most variable}
(42--\SI{71}{cm}); accuracy \SI{\rsMETRICACC}{cm} vs completeness \SI{\rsMETRICCOMP}{cm}: the surfaces are all there but
in the wrong place (\S\ref{sec:metric}). Kitchen 45261556, where mono$+$oracle is second best (\rsORACLEf), is where the
metric model is worst (\rsMETRICf)---right shape, wrong scale. This is the baseline of \S\ref{sec:main}: fine-tuning this
same model on ARKitScenes with a LiDAR teacher removes most of that error.
\end{enumerate}

\begin{figure*}[t]
\centering
\includegraphics[width=\textwidth]{exp1_depth_source_topdown_grid}
\caption{Exp.~1 top-down error maps (rows $=$ scenes, columns $=$ runs, colour $=$ distance to reference,
0--\SI{10}{cm}). The gt/lidar columns are almost entirely blue ($<$\SI{2}{cm}); oracle and sparse show the same
yellow--red pattern at object edges and far walls; per\_scene has a red blob in the middle of every room (phantom surfaces from wrongly scaled
frames); metric is red almost everywhere.}
\label{fig:exp1}
\end{figure*}

\subsection{Context, Exp.~2: voxel size (GT depth)}
\begin{table}[h]
\centering\small
\caption{\code{depth.source: gt}, \code{sdf\_trunc} $=3\times$ voxel.}
\begin{adjustbox}{max width=\columnwidth}
\begin{tabular}{@{}lrrrrrrr@{}}
\toprule
Voxel & Chamfer & Acc. & Compl. & F@2 & F@5 & F@10 & Fusion (s) \\
\midrule
\SI{2}{cm} & $1.18 \pm 0.19$ & 1.05 & 1.31 & 0.92 & 0.99 & 1.00 & 2.5 \\
\textbf{\SI{4}{cm}} & $2.06 \pm 0.54$ & 1.29 & 2.83 & 0.86 & 0.97 & 0.99 & 1.5 \\
\SI{8}{cm} & $2.76 \pm 0.85$ & 1.74 & 3.79 & 0.76 & 0.93 & 0.97 & 1.3 \\
\bottomrule
\end{tabular}
\end{adjustbox}
\end{table}
Going from 2 to \SI{4}{cm} costs \SI{0.9}{cm} Chamfer, mostly completeness (1.3$\rightarrow$2.8: coarser corners and
edges), not accuracy (1.05$\rightarrow$1.29). From 4 to \SI{8}{cm} costs another \SI{0.7}{cm} and \fscore{} drops
0.97$\rightarrow$0.93---at \SI{8}{cm} discretisation exceeds LiDAR's own error. Fusion takes 1--\SI{3}{s} per scan at
every size (CPU, Open3D), so voxel size does not drive cost; inference does (\S\ref{sec:business}). \textbf{\SI{4}{cm}
is used throughout}: its discretisation (\SI{2.1}{cm}) is smaller than the depth-source effects being measured (LiDAR
$+1$, mono $+4$, per\_scene $+20$) and the mesh stays small enough for a phone.

\subsection{Context, Exp.~3: frame stride---how many frames per second to capture}\label{sec:stride}
Faro provides frames at $\sim$4.1~fps (2{,}352 frames / \SI{577}{s}); strides 1/5/10/20 therefore correspond to
$\approx$ 4 / 0.8 / 0.4 / 0.2~fps (391 / 78 / 40 / 20 frames per scan on average). Three rows per stride separate the
causes: GT control (coverage only), mono$+$oracle ($+$ model shape), mono$+$per\_scene ($+$ scale).
\begin{table*}[t]
\centering\small
\caption{Chamfer / accuracy / completeness (cm) and \fscore{} vs.\ stride.}
\begin{adjustbox}{max width=\textwidth}
\begin{tabular}{@{}lrrrrrrrr@{}}
\toprule
Stride ($\approx$fps) & Frames & GT: Ch / acc / comp & $F$ & mono oracle: Ch / acc / comp & $F$ & mono per\_scene: Ch / acc / comp & $F$ & mono time (s) \\
\midrule
1 (4) & \rsSTRGTONEN & \rsSTRGTONECM / \textbf{\rsSTRGTONEACC} / \rsSTRGTONECOMP & \rsSTRGTONEF & \rsSTRORACLEONECM / \rsSTRORACLEONEACC / \rsSTRORACLEONECOMP & \rsSTRORACLEONEF & \rsSTRSCENEONECM / \rsSTRSCENEONEACC / \rsSTRSCENEONECOMP & \rsSTRSCENEONEF & \rsSTRSCENEONESCAN \\
5 (0.8) & \rsSTRGTFIVEN & \rsSTRGTFIVECM / \textbf{\rsSTRGTFIVEACC} / \rsSTRGTFIVECOMP & \rsSTRGTFIVEF & \rsSTRORACLEFIVECM / \rsSTRORACLEFIVEACC / \rsSTRORACLEFIVECOMP & \rsSTRORACLEFIVEF & \rsSTRSCENEFIVECM / \rsSTRSCENEFIVEACC / \rsSTRSCENEFIVECOMP & \rsSTRSCENEFIVEF & \rsSTRSCENEFIVESCAN \\
10 (0.4) & \rsSTRGTTENN & \rsSTRGTTENCM / \textbf{\rsSTRGTTENACC} / \rsSTRGTTENCOMP & \rsSTRGTTENF & \rsSTRORACLETENCM / \rsSTRORACLETENACC / \rsSTRORACLETENCOMP & \rsSTRORACLETENF & \rsSTRSCENETENCM / \rsSTRSCENETENACC / \rsSTRSCENETENCOMP & \rsSTRSCENETENF & \rsSTRSCENETENSCAN \\
20 (0.2) & \rsSTRGTTWENTYN & \rsSTRGTTWENTYCM / \textbf{\rsSTRGTTWENTYACC} / \rsSTRGTTWENTYCOMP & \rsSTRGTTWENTYF & \rsSTRORACLETWENTYCM / \rsSTRORACLETWENTYACC / \rsSTRORACLETWENTYCOMP & \rsSTRORACLETWENTYF & \rsSTRSCENETWENTYCM / \rsSTRSCENETWENTYACC / \rsSTRSCENETWENTYCOMP & \rsSTRSCENETWENTYF & \rsSTRSCENETWENTYSCAN \\
\bottomrule
\end{tabular}
\end{adjustbox}
\end{table*}
\begin{itemize}
\item \textbf{Control:} GT accuracy is constant at 1.2--\SI{1.3}{cm} at every stride; the growing Chamfer
(\rsSTRGTONECM$\rightarrow$\SI{\rsSTRGTTWENTYCM}{cm}) is pure completeness. On this dataset stride measures \textbf{coverage}: a room swept once
develops holes below 1~fps.
\item \textbf{mono$+$oracle tracks the control plus a constant \SIrange{2.7}{3.3}{cm}}, and its accuracy
\emph{improves} with stride (\rsSTRORACLEONEACC$\rightarrow$\rsSTRORACLETWENTYACC): fewer frames means fewer
overlapping surfaces from different views; 2D AbsRel is \rsSTRORACLEONEABSREL{} at every stride, confirming the model
did not change.
\item \textbf{mono$+$per\_scene is flat at $\sim$\SI{23}{cm} up to stride 10}: at stride 5 it saves $5\times$ the time
(\rsSTRSCENEONESCAN$\rightarrow$\SI{\rsSTRSCENEFIVESCAN}{s}) at the same Chamfer (\rsSTRSCENEONECM{} vs.\
\rsSTRSCENEFIVECM), because the error is already dominated by scale. At stride 20
\code{per\_scene} fits on 10 of 20 frames (nearly a half-scene oracle), so that cell should not be over-interpreted.
\item \textbf{For the business:} with good depth (LiDAR) capture $\geq 1$~fps for a complete room (\fscore{} $>0.9$);
with mono$+$per\_scene, 0.8~fps gives the same result as 4~fps, $5\times$ faster---quality is capped by scale, not by
frame count.
\end{itemize}
At stride 20 every row has the same holes (areas the camera passed quickly): holes are a property of the scan, not of
the depth source.

\subsection{Context, Exp.~4: model size (oracle scale)}
\begin{table}[h]
\centering\small
\caption{\code{oracle\_frame} for every model, \SI{4}{cm}, every frame---pure ``shape,'' excluding the scale problem of \S5.1.}
\begin{adjustbox}{max width=\columnwidth}
\begin{tabular}{@{}lrrrrrrr@{}}
\toprule
Model & Params & Chamfer & F@5 & AbsRel & $\delta_1$ & ms/frame (MPS) & s/scan \\
\midrule
DA-v2 Large & 335M & $\mathbf{\rsDALARGECM \pm \rsDALARGESD}$ & \rsDALARGEF & \rsDALARGEABSREL & \rsDALARGEDELTAONE & 871 & \rsDALARGESCAN \\
DA-v2 Small & 25M & $\mathbf{\rsDASMALLCM \pm \rsDASMALLSD}$ & \rsDASMALLF & \rsDASMALLABSREL & \rsDASMALLDELTAONE & 121 & \rsDASMALLSCAN \\
MiDaS v2.1 small & 21M & $\mathbf{\rsMIDASCM \pm \rsMIDASSD}$ & \rsMIDASF & \rsMIDASABSREL & \rsMIDASDELTAONE & 27 & \rsMIDASSCAN \\
\bottomrule
\end{tabular}
\end{adjustbox}
\end{table}
L$\rightarrow$S: $13\times$ smaller, $6\times$ faster, $+$\SI{1.0}{cm} ($F$ \rsDALARGEF$\rightarrow$\rsDASMALLF);
every scene moves the same way except the bathroom, where Small is slightly \emph{better} (7.5$\rightarrow$7.1).
S$\rightarrow$MiDaS: similar parameter count, a further $3\times$ faster, but $2.1\times$ the Chamfer and $F$
\rsMIDASF---a matter of model generation (DA-v2 was trained with 62M unlabelled images), not size. At oracle scale
DA-v2 S (\SI{\rsDASMALLCM}{cm}) is still \SI{3.3}{cm} behind LiDAR (\rsLIDARCM).
Model size is not the main bottleneck: Table~\ref{tab:exp1} shows scale (oracle$\rightarrow$per\_scene,
$+$\SI{17.6}{cm}) costs $17\times$ more than model size (L$\rightarrow$S, $+$\SI{1.0}{cm}).

\subsection{Context, Exp.~5: does LiDAR confidence help?}\label{sec:exp5}
ARKit ships a $\{0,1,2\}$ confidence map with every LiDAR frame. Masking or weighting by it is the cheapest possible
improvement to the ``iPhone~Pro today'' row.
\begin{table}[h]
\centering\small
\caption{\code{lidar} row with confidence masking / integer TSDF weights (\SI{4}{cm}, every frame).}
\begin{adjustbox}{max width=\columnwidth}
\begin{tabular}{@{}lrrrr@{}}
\toprule
Variant & Chamfer & Acc. & Compl. & F@5 \\
\midrule
all pixels, weight 1 (Exp.~1) & $3.04 \pm 0.57$ & 2.23 & 3.85 & 0.93 \\
mask confidence $\geq 1$ & \rsCONFONECM & \rsCONFONEACC & \rsCONFONECOMP & \rsCONFONEF \\
mask confidence $= 2$ & \rsCONFTWOCM & \rsCONFTWOACC & \rsCONFTWOCOMP & \rsCONFTWOF \\
weights $[0,1,2]$ & \rsWZOTCM & \rsWZOTACC & \rsWZOTCOMP & \rsWZOTF \\
weights $[1,2,4]$ & \rsWOTFCM & \rsWOTFACC & \rsWOTFCOMP & \rsWOTFF \\
\bottomrule
\end{tabular}
\end{adjustbox}
\end{table}
Nothing moves: dropping confidence-0 pixels or weighting by tier changes Chamfer by $\leq$\SI{0.03}{cm}; keeping only
confidence-2 pixels improves accuracy slightly (2.23$\rightarrow$2.08) but loses completeness (3.85$\rightarrow$5.23,
worst in the closet-heavy bedroom 47429736: 3.1$\rightarrow$5.8). The \SI{1}{cm} LiDAR$-$laser gap is not a
low-confidence-pixel problem; it is the sensor's resolution and its 2--\SI{4}{cm} noise everywhere.

\subsection{Numbers carried into \S6--7}
\begin{table}[h]
\centering\small
\begin{tabularx}{\columnwidth}{@{}lX@{}}
\toprule
Question & Number \\
\midrule
\textbf{R3 vs.\ pretr.\ (3D, 6 rooms)} & \textbf{\SI{\rsMETRICCM}{cm} $\rightarrow$ \SI{\rsFTLIDARALLCM}{cm}, \fscore{} \rsMETRICF{} $\rightarrow$ \rsFTLIDARALLF, every room better} \\
\textbf{R3 vs.\ pretr.\ (2D, 6 rooms)} & \textbf{AbsRel \rsMETRICSIXABSREL{} $\rightarrow$ \rsFTLIDARALLABSREL, $\delta_1$ \rsMETRICSIXDELTAONE{} $\rightarrow$ \rsFTLIDARALLDELTAONE} \\
\textbf{R3 vs.\ pretr.\ (2D, \rsVN{} rooms)} & \textbf{AbsRel \rsVMETRICABSREL{} $\rightarrow$ \rsVFTALLABSREL, better in \rsVWINS/\rsVN} \\
Faro vs.\ LiDAR teacher & \SI{\rsFTFAROCM}{cm} / \SI{\rsFTLIDARCM}{cm} / \SI{\rsFTLIDARALLCM}{cm} (R1/R2/R3): not measurably different \\
pipeline loss at \SI{4}{cm} & \SI{2.1 \pm 0.5}{cm} \\
iPad LiDAR vs laser & $+$\SI{1.0 \pm 0.3}{cm}, \fscore{} 0.93 (confidence weighting: no change) \\
mono ceiling (scale known per frame) & $+$\SI{\rsORACLEGAPLIDARCM \pm \rsORACLEGAPLIDARSD}{cm} from LiDAR, \fscore{} \rsORACLEF \\
mono with 200 sparse points/frame & \SI{\rsSPARSECM}{cm}, \fscore{} \rsSPARSEF{} (\SI{\rsSPARSEGAP}{cm} from oracle) \\
price of ``one scale per scene'' & $4.5\times$ oracle $\rightarrow$ \SI{\rsSCENECM}{cm}, \fscore{} \rsSCENEF \\
intrinsics-conditioned model (Depth Pro) & \SI{\rsDEPTHPROCM}{cm}, \fscore{} \rsDEPTHPROF{} --- intrinsics do not replace in-domain labels \\
voxel 2$\rightarrow$4$\rightarrow$\SI{8}{cm} & 1.2 $\rightarrow$ 2.1 $\rightarrow$ \SI{2.8}{cm} \\
fps needed (LiDAR/GT) & $\geq 0.8$~fps for \fscore{} $\geq 0.9$; 0.4~fps leaves 0.83 \\
L$\rightarrow$S$\rightarrow$MiDaS & \rsDALARGECM $\rightarrow$ \rsDASMALLCM $\rightarrow$ \SI{\rsMIDASCM}{cm} at 871 $\rightarrow$ 121 $\rightarrow$ 27 ms/frame \\
\bottomrule
\end{tabularx}
\end{table}
